\documentclass[../main.tex]{subfiles}

\begin{document}

\chapter{Contexte général du projet}
\label{chap:contexte}

\section{Atlantic Re : un réassureur international à ancrage marocain}

\subsection{Le Groupe CDG, actionnaire de référence}

La Caisse de Dépôt et de Gestion (CDG) est l'institution financière publique de
référence au Maroc, fondée en 1959. Institution de long terme, elle joue un rôle
structurant dans le financement des grandes infrastructures nationales et porte un
portefeuille de filiales couvrant la finance, l'assurance et l'investissement. Elle
détient plus de 94\,\% du capital d'Atlantic Re, dont elle constitue l'actionnaire
stratégique et le socle de crédibilité institutionnelle.

\subsection{Profil et périmètre mondial d'Atlantic Re}

Atlantic Re trouve ses racines dans la \textbf{Société Centrale de Réassurance
(SCR)}, première compagnie de réassurance du Maroc, créée en \textbf{1960}. Pendant
plus de six décennies, la SCR a occupé une position prépondérante sur le marché
marocain avant d'étendre ses activités à l'international. En \textbf{2024}, le
rebranding en \textbf{Atlantic Re} traduit l'envergure désormais multinationale de
la structure, dont les opérations couvrent l'Afrique, le Moyen-Orient et l'Asie.
Elle souscrit des risques en branches Vie et Non-Vie, couvrant des lignes d'affaires
aussi diverses que l'incendie, le transport, l'automobile, la responsabilité civile,
les risques techniques, le décès et la prévoyance.

Au terme de l'exercice 2025, Atlantic Re enregistre une progression soutenue de son
activité, confirmant la dynamique engagée dans le cadre de son plan stratégique
\textbf{Reach2030}. Le tableau~\ref{tab:kpi_atlanticre} synthétise les indicateurs
clés de performance publiés pour cet exercice.

\begin{table}[H]
\centering
\caption{Indicateurs clés de performance d'Atlantic Re --- Exercice 2025}
\label{tab:kpi_atlanticre}
\small
\begin{tabularx}{\textwidth}{@{}p{5cm}p{3cm}X@{}}
\toprule
\textbf{Indicateur} & \textbf{Valeur 2025} & \textbf{Commentaire} \\
\midrule
Chiffre d'affaires            & \textbf{4~039 MDH}  & En dépassement des projections annuelles \\
Résultat net                  & \textbf{419 MDH}    & Impact positif du plan Reach2030 \\
Ratio combiné                 & \textbf{92,3\,\%}   & Sinistralité favorable ; renforcement de la prudence \\
Ratio de solvabilité S2       & \textbf{215\,\%}    & En hausse de 9 points vs exercice précédent \\
ROE (Return On Equity)        & \textbf{13,4\,\%}   & Rendement annualisé des capitaux propres \\
Présence géographique         & \textbf{+60 pays}   & Couverture multinationale \\
Portefeuille clients          & \textbf{+400 clients} & Confiance renouvelée des partenaires \\
\bottomrule
\end{tabularx}
\end{table}

Ces performances s'inscrivent dans la continuité de la dynamique de développement
engagée par Atlantic Re, portée par le renforcement de sa présence sur les marchés
internationaux, la consolidation de sa politique de souscription et la valorisation
de ses placements financiers. Le ratio de solvabilité S2 à 215\,\% --- en hausse de
9 points --- témoigne du renforcement des fonds propres économiques et d'une maîtrise
accrue des capitaux sous risques, deux prérequis essentiels à toute stratégie
d'expansion internationale.

C'est dans ce contexte de solidité financière et d'ambition continentale qu'Atlantic
Re confirme son positionnement en tant que réassureur de référence au service du
marché marocain et des marchés régionaux, tout en poursuivant une trajectoire de
croissance durable fondée sur la performance et la création de valeur.

Sur la scène africaine, Atlantic Re évolue dans un environnement concurrentiel dense
et stratifié. Le paysage réassurantiel continental regroupe des acteurs de natures
très diverses : des \textbf{réassureurs panafricains institutionnels} tels qu'Africa
Re (premier réassureur panafricain, dont Atlantic Re est actionnaire), ZEP-Re
(spécialisé en Afrique orientale et australe) et CICA-Re (réassureur institutionnel
de la zone CIMA) ; des \textbf{réassureurs internationaux de premier rang} (Munich
Re, Swiss Re, Hannover Re, Scor) qui opèrent depuis leurs sièges européens en
souscrivant des portefeuilles africains en ligne ou via des bureaux régionaux ; des
\textbf{réassureurs régionaux et locaux} implantés dans les zones anglophone, lusophone
ou CIMA, qui capitalisent sur leur proximité réglementaire et leurs relations de
longue date avec les cédantes locales ; enfin, des \textbf{courtiers de réassurance
de dimension mondiale} (Aon, Guy Carpenter, Willis Re) qui, par leur capacité de
placement et leur accès à la capacité internationale, exercent une influence croissante
sur la structuration des traités en Afrique. Atlantic Re capitalise dans cet
écosystème sur un ancrage marocain conférant crédibilité et neutralité continentale,
une expertise des marchés francophones construite sur six décennies, et le soutien
institutionnel du Groupe CDG.

\subsection{Le Pôle Business Développement : cadre d'accueil du stage}

Le stage se déroule au sein du \textbf{Pôle Business Développement} de la Direction
Technique et Développement d'Atlantic Re, sur une durée de \textbf{quatre mois}
(février--juin 2026). Ce pôle est en charge de l'identification et de l'analyse des
opportunités de développement commercial sur les marchés existants et potentiels. Il
produit les études de marché qui alimentent les décisions de souscription de la
direction générale. C'est dans ce cadre que s'inscrit la commande initiale du stage :
fournir un modèle quantitatif de priorisation des marchés africains à l'horizon 2030.

\section{Contexte stratégique et objectif initial du PFE}

\subsection{Le marché africain : une opportunité structurelle}

Le marché africain de l'assurance Vie et Non-Vie est en transformation profonde. La
croissance démographique, l'urbanisation accélérée, les grands projets d'infrastructure
et la montée en puissance des risques climatiques créent une demande structurelle
de couverture assurantielle que les capacités locales peinent à satisfaire seules.
Cette situation génère un besoin de réassurance que des acteurs comme Atlantic Re
sont en position de capter, à condition de disposer d'une information fiable et
d'un outil de priorisation stratégique des marchés \citep{Guerard2021}.

C'est précisément dans ce cadre que s'inscrit le plan stratégique \textbf{Reach2030}
d'Atlantic Re, qui fixe l'expansion africaine comme axe de croissance prioritaire à
l'horizon 2030. Avec un ratio de solvabilité S2 établi à 215\,\% et un ROE de
13,4\,\% pour l'exercice 2025, Atlantic Re dispose d'une assise financière solide
pour financer cette expansion, tout en maintenant une discipline technique illustrée
par un ratio combiné de 92,3\,\%. La couverture actuelle de plus de 60 pays et un
portefeuille dépassant 400 clients témoignent de la capacité opérationnelle du
réassureur à opérer dans des environnements réglementaires et économiques variés.

\subsection{La problématique et l'élargissement du périmètre}

Atlantic Re opère sur plusieurs marchés africains mais sans modèle formalisé
permettant de comparer objectivement leurs profils d'attractivité, de rentabilité
et de risque. La problématique s'articule donc autour de la sélection optimale
des marchés africains offrant le meilleur potentiel de développement selon une
démarche multicritère quantifiée. Cependant, l'absence d'outils de gouvernance
pour les données internes du portefeuille a nécessité un élargissement du
périmètre de la mission. Ainsi, le projet intègre désormais deux axes
complémentaires et indissociables : la \textbf{digitalisation} (développement
d'un écosystème digital pour la fiabilisation et la consultation de la donnée interne)
servant de socle préalable incontournable à la \textbf{modélisation} (pipeline
data science et scoring SCAR des marchés africains).

% ============================================================
%  SECTION 1.3 — OBJECTIFS
% ============================================================
\section{Objectifs du projet}

Le projet poursuit les objectifs opérationnels suivants, chacun répondant à un besoin identifié lors de l'analyse de terrain et validé par les parties prenantes :

\begin{itemize}[leftmargin=1.5cm, itemsep=8pt]

  \item \textbf{Auditer et assainir la base de données interne d'Atlantic Re} afin d'établir un diagnostic complet de la qualité du portefeuille de réassurance, d'identifier les incohérences structurelles et de produire un référentiel de données normalisé servant de socle fiable aux analyses ultérieures.

  \item \textbf{Concevoir et développer un écosystème digital de gouvernance de la donnée interne} offrant aux équipes du Pôle Business Développement un portail structuré de consultation, de filtrage multicritère et d'export des indicateurs clés de performance du portefeuille, avec des visualisations interactives, des modules d'analyse métier dédiés et un système de gestion des accès sécurisé par profils.

  \item \textbf{Constituer un panel de données externes couvrant les pays du continent africain sur dix exercices (2015--2024)} en collectant et harmonisant des données issues de sources institutionnelles reconnues, et en garantissant la complétude et la traçabilité de l'ensemble du jeu de données produit.

  \item \textbf{Développer des tableaux de bord cartographiques interactifs de l'espace africain} permettant de visualiser, comparer et analyser l'attractivité des marchés selon quatre dimensions complémentaires : marché assurance Non-Vie, marché assurance Vie, environnement macroéconomique et gouvernance institutionnelle.

  \item \textbf{Construire et calibrer le modèle de scoring multicritère SCAR} (Scoring par Corrélation Ajustée et Robustesse) pour produire un classement opérationnel des marchés africains, assorti de recommandations stratégiques directement exploitables par la direction dans le cadre du plan Reach2030 --- \textit{cette composante est en cours de développement au moment de la rédaction du rapport}.

  \item \textbf{Mettre en place un croisement entre les données internes et externes} permettant d'évaluer l'exposition relative d'Atlantic Re sur chaque marché africain, d'identifier les marchés sous-exploités et de quantifier le potentiel de croissance du réassureur sur le continent.

\end{itemize}

% ============================================================
%  SECTION 1.4 — MÉTHODES DE TRAVAIL
% ============================================================
\section{Méthodes de travail}

\subsection{Approche Agile adaptée}

Le projet a été conduit selon une approche \textbf{Agile adaptée}, inspirée du cadre Scrum mais ajustée au contexte d'un PFE de quatre mois réalisé par un binôme. Cette méthodologie repose sur un découpage du travail en \textbf{sprints de deux à quatre semaines}, chacun produisant un livrable partiel démontrable et validable par l'encadrant entreprise. L'adoption de cette démarche itérative et incrémentale se justifie par trois facteurs principaux : la nature exploratoire du projet --- dont le périmètre a évolué au terme de la phase d'immersion ---, la nécessité de livrables intermédiaires pour maintenir la confiance du commanditaire, et la complémentarité des deux axes du projet qui nécessitent un parallélisme de tâches maîtrisé entre les deux membres du binôme.

\subsection{Rituels et outils de gestion}

Les rituels suivants ont structuré le suivi du projet :

\begin{itemize}[leftmargin=1.5cm, itemsep=4pt]
  \item \textbf{Point hebdomadaire avec l'encadrant entreprise} : revue des livrables de la semaine, validation des orientations, arbitrage des priorités et synchronisation avec les contraintes métier du Pôle Business Développement.
  \item \textbf{Point hebdomadaire avec l'encadrante académique} : revue des avancements méthodologiques, validation de la rigueur scientifique des choix réalisés et ajustement du plan de rédaction du rapport.
  \item \textbf{Répartition des tâches au sein du binôme} : chaque sprint fait l'objet d'une répartition explicite entre les deux membres, avec des zones de responsabilité distinctes et des points de synchronisation réguliers.
  \item \textbf{Contrôle de version du code source} : l'ensemble des productions numériques du projet est versionné sur un dépôt privé avec des historiques traçables et des revues de code croisées entre les deux membres.
  \item \textbf{Tableau de suivi des activités} : les livrables et les tâches associées sont suivis via un tableau de bord Kanban maintenu tout au long du projet, permettant de visualiser l'avancement en temps réel.
\end{itemize}

\subsection{Cycle de développement}

Chaque sprint suit un cycle en quatre temps : (1)~\textbf{planification} --- définition des objectifs du sprint et des critères d'acceptation ; (2)~\textbf{réalisation} --- développement itératif des composants avec intégration régulière des productions des deux membres ; (3)~\textbf{revue} --- démonstration des livrables au commanditaire et recueil du retour métier ; (4)~\textbf{rétrospective} --- identification des ajustements de processus pour le sprint suivant. Ce cycle garantit que les livrables évoluent en fonction du retour terrain et que les risques sont détectés au plus tôt.

\clearpage
% ============================================================
%  SECTION 1.5 — PLANNING DU PROJET
% ============================================================
\section{Planning du projet}

Le tableau~\ref{tab:planning} présente le planning détaillé du projet, organisé en huit sprints couvrant la totalité de la durée du stage (février à juin 2026). Ce planning reflète la trajectoire réelle du projet, incluant l'élargissement du périmètre intervenu au terme du Sprint~0 et la décomposition fine des activités en livrables intermédiaires validables.

\begin{table}[H]
\centering
\caption{Planning détaillé du projet par sprints}
\label{tab:planning}
\small
\renewcommand{\arraystretch}{1.35}
\begin{tabularx}{\textwidth}{@{}p{1.7cm}p{2.9cm}X@{}}
\toprule
\textbf{Sprint} & \textbf{Période} & \textbf{Activités et livrables} \\
\midrule

\rowcolor{phase1}
\textbf{Sprint 0} \newline \textit{Immersion \& Cadrage}
& 03 fév. -- 21 fév. \newline (3 semaines)
& \textbullet\ Découverte du métier de réassureur et du Pôle Business Développement \newline
  \textbullet\ Étude du plan stratégique Reach2030 et des enjeux de l'expansion africaine \newline
  \textbullet\ Audit de la base de données interne : diagnostic qualité et identification des problèmes structurels \newline
  \textbullet\ Détection du besoin de digitalisation et proposition d'élargissement du périmètre \newline
  \textbf{Livrable :} rapport d'audit interne + note de cadrage validée \\

\midrule
\rowcolor{phase2}
\textbf{Sprint 1} \newline \textit{Analyse des besoins \& Conception}
& 24 fév. -- 07 mars \newline (2 semaines)
& \textbullet\ Formalisation des besoins explicites et cachés selon la méthode MoSCoW \newline
  \textbullet\ Cartographie des acteurs et des flux de données internes \newline
  \textbullet\ Conception de l'architecture logique du système et du modèle de données \newline
  \textbullet\ Élaboration des diagrammes UML (cas d'utilisation, classes, séquence) \newline
  \textbf{Livrable :} cahier des charges validé par les parties prenantes \\

\midrule
\rowcolor{phase3}
\textbf{Sprint 2} \newline \textit{Architecture \& Infrastructure}
& 10 mars -- 21 mars \newline (2 semaines)
& \textbullet\ Définition de l'architecture technique et des composants du système \newline
  \textbullet\ Mise en place de l'environnement de développement et du contrôle de version \newline
  \textbullet\ Conception du modèle de sécurité et de gestion des accès par profils \newline
  \textbullet\ Développement du service de chargement et de normalisation des données internes \newline
  \textbf{Livrable :} infrastructure opérationnelle et premier écran fonctionnel \\

\midrule
\rowcolor{phase4}
\textbf{Sprint 3} \newline \textit{Portail --- Tableau de bord \& Filtrage}
& 24 mars -- 11 avr. \newline (3 semaines)
& \textbullet\ Développement du tableau de bord principal avec indicateurs clés du portefeuille \newline
  \textbullet\ Mise en place du système de filtrage multicritère (pays, branche, exercice, cédante\ldots) \newline
  \textbullet\ Développement des visualisations interactives associées aux filtres \newline
  \textbullet\ Gestion des utilisateurs, des rôles et journal d'activité \newline
  \textbullet\ Export des données filtrées aux formats standards \newline
  \textbf{Livrable :} tableau de bord interactif opérationnel et validé \\

\midrule
\rowcolor{phase5}
\textbf{Sprint 4} \newline \textit{Portail --- Modules métier avancés}
& 14 avr. -- 02 mai \newline (3 semaines)
& \textbullet\ Module d'analyse par cédante avec réconciliation intelligente des dénominations \newline
  \textbullet\ Module d'analyse par courtier avec indicateurs de performance \newline
  \textbullet\ Module de comparaison directe entre marchés \newline
  \textbullet\ Module de pilotage des parts cibles et de saturation des affaires facultatives \newline
  \textbullet\ Module de rétrocession avec analyse des partenaires \newline
  \textbf{Livrable :} portail de gouvernance complet (MVP Axe 1) \\

\midrule
\rowcolor{phase6}
\textbf{Sprint 5} \newline \textit{Données externes \& Cartographies}
& 05 mai -- 30 mai \newline (4 semaines)
& \textbullet\ Collecte, harmonisation et structuration des données de 34 marchés africains (2015--2024) \newline
  \textbullet\ Traitement des valeurs manquantes par pipeline hiérarchique en trois phases \newline
  \textbullet\ Développement des quatre modules cartographiques thématiques (Non-Vie, Vie, Macro, Gouvernance) \newline
  \textbullet\ Module d'analyse par pays et module de comparaison côte-à-côte \newline
  \textbullet\ Croisement des données internes et externes par marché \newline
  \textbf{Livrable :} panel africain complet et visualisations cartographiques validées \\

\midrule
\rowcolor{phase7}
\textbf{Sprint 6} \newline \textit{Modélisation SCAR}
& 02 juin -- 13 juin \newline (2 semaines)
& \textbullet\ Construction et calibration du modèle de scoring multicritère SCAR \newline
  \textbullet\ Validation statistique des résultats (procédure Leave-One-Country-Out) \newline
  \textbullet\ Production du classement opérationnel des marchés africains \newline
  \textbullet\ Génération des recommandations stratégiques par marché \newline
  \textit{Note : cette phase est en cours de réalisation au moment de la rédaction du rapport} \newline
  \textbf{Livrable :} classement SCAR et note de recommandations stratégiques \\

\midrule
\rowcolor{phase8}
\textbf{Sprint 7} \newline \textit{Intégration \& Clôture}
& 16 juin -- 27 juin \newline (2 semaines)
& \textbullet\ Tests d'intégration et corrections sur l'ensemble du système \newline
  \textbullet\ Validation métier des résultats avec l'encadrant entreprise \newline
  \textbullet\ Finalisation de la documentation et du rapport de PFE \newline
  \textbullet\ Préparation et répétition de la soutenance \newline
  \textbf{Livrable :} rapport final, application livrée et présentation soutenance \\

\bottomrule
\end{tabularx}
\end{table}

La figure~\ref{fig:gantt} présente le diagramme de Gantt du projet.

\begin{figure}[H]
\centering
\caption{Diagramme de Gantt --- Planning du projet (Février -- Juin 2026)}
\label{fig:gantt}
\begin{tikzpicture}[x=0.545cm, y=0.78cm]

  % ── Unités : 1 unité x = 1 semaine ; 20 semaines au total
  % ── Colonne libellés : x = -8 à 0  |  Timeline : x = 0 à 20
  % ── Mois : Fév(0-4) Mars(4-8) Avr(8-12) Mai(12-16) Juin(16-20)

  % ─── Fonds alternés par mois ────────────────────────────────
  \fill[gray!10] (0,0) rectangle (4,9);
  \fill[white]   (4,0) rectangle (8,9);
  \fill[gray!10] (8,0) rectangle (12,9);
  \fill[white]   (12,0) rectangle (16,9);
  \fill[gray!10] (16,0) rectangle (20,9);

  % ─── En-tête ────────────────────────────────────────────────
  \fill[gray!65] (-9.5,8) rectangle (20,9);
  \node[font=\small\bfseries, text=white, anchor=west] at (-9.3,8.5)
    {Phase / Sprint};
  \foreach \x/\lbl in {0/Février, 4/Mars, 8/Avril, 12/Mai, 16/Juin} {
    \node[font=\small\bfseries, text=white] at (\x+2,8.5) {\lbl};
  }

  % ─── Séparateurs verticaux mois (en-tête) ───────────────────
  \foreach \x in {0,4,8,12,16} {
    \draw[gray!40] (\x,8) -- (\x,9);
  }

  % ─── Grille horizontale (lignes de rangées) ─────────────────
  \foreach \y in {0,1,...,8} {
    \draw[gray!35, thin] (-9.5,\y) -- (20,\y);
  }

  % ─── Grille verticale semaines (très fine) ──────────────────
  \foreach \w in {1,2,3,5,6,7,9,10,11,13,14,15,17,18,19} {
    \draw[gray!18, very thin] (\w,0) -- (\w,8);
  }
  \foreach \x in {0,4,8,12,16,20} {
    \draw[gray!45] (\x,0) -- (\x,8);
  }

  % ─── Séparateur colonne libellés / timeline ─────────────────
  \draw[gray!55, semithick] (0,0) -- (0,9);

  % ─── Bordure extérieure ─────────────────────────────────────
  \draw[gray!55, semithick] (-9.5,0) rectangle (20,9);

  % ─── Libellés des sprints (colonne gauche) ──────────────────
  \node[font=\scriptsize\bfseries, anchor=west] at (-9.3,7.5)
    {S0 \textemdash\ Immersion \& Cadrage};
  \node[font=\scriptsize\bfseries, anchor=west] at (-9.3,6.5)
    {S1 \textemdash\ Analyse \& Conception};
  \node[font=\scriptsize\bfseries, anchor=west] at (-9.3,5.5)
    {S2 \textemdash\ Architecture \& Infrastructure};
  \node[font=\scriptsize\bfseries, anchor=west] at (-9.3,4.5)
    {S3 \textemdash\ Portail : Tableau de bord};
  \node[font=\scriptsize\bfseries, anchor=west] at (-9.3,3.5)
    {S4 \textemdash\ Portail : Modules métier};
  \node[font=\scriptsize\bfseries, anchor=west] at (-9.3,2.5)
    {S5 \textemdash\ Données ext. \& Cartographies};
  \node[font=\scriptsize\bfseries, anchor=west] at (-9.3,1.5)
    {S6 \textemdash\ Modélisation SCAR};
  \node[font=\scriptsize\bfseries, anchor=west] at (-9.3,0.5)
    {S7 \textemdash\ Intégration \& Clôture};

  % ─── Barres de sprints ──────────────────────────────────────
  % S0 : sem. 0–3
  \fill[phase1, draw=phase1!55!black, rounded corners=2pt]
    (0.04,7.12) rectangle (3,7.88);
  \node[font=\tiny\bfseries, text=gray!20!black] at (1.5,7.5) {3 sem.};

  % S1 : sem. 3–5
  \fill[phase2, draw=phase2!55!black, rounded corners=2pt]
    (3,6.12) rectangle (5,6.88);
  \node[font=\tiny\bfseries, text=gray!20!black] at (4,6.5) {2 sem.};

  % S2 : sem. 5–7
  \fill[phase3, draw=phase3!55!black, rounded corners=2pt]
    (5,5.12) rectangle (7,5.88);
  \node[font=\tiny\bfseries, text=gray!20!black] at (6,5.5) {2 sem.};

  % S3 : sem. 7–10
  \fill[phase4, draw=phase4!55!black, rounded corners=2pt]
    (7,4.12) rectangle (10,4.88);
  \node[font=\tiny\bfseries, text=gray!20!black] at (8.5,4.5) {3 sem.};

  % S4 : sem. 10–13
  \fill[phase5, draw=phase5!55!black, rounded corners=2pt]
    (10,3.12) rectangle (13,3.88);
  \node[font=\tiny\bfseries, text=gray!20!black] at (11.5,3.5) {3 sem.};

  % S5 : sem. 13–17
  \fill[phase6, draw=phase6!55!black, rounded corners=2pt]
    (13,2.12) rectangle (17,2.88);
  \node[font=\tiny\bfseries, text=gray!20!black] at (15,2.5) {4 sem.};

  % S6 : sem. 17–19
  \fill[phase7, draw=phase7!55!black, rounded corners=2pt]
    (17,1.12) rectangle (19,1.88);
  \node[font=\tiny\bfseries, text=gray!20!black] at (18,1.5) {2 sem.};

  % S7 : sem. 18–20
  \fill[phase8, draw=gray!55, rounded corners=2pt]
    (18,0.12) rectangle (20,0.88);
  \node[font=\tiny\bfseries, text=gray!30!black] at (19,0.5) {2 sem.};

  % ─── Jalons ─────────────────────────────────────────────────
  % ◆ Élargissement du périmètre (fin S0 — sem. 3)
  \node[diamond, fill=red!72, draw=red!80!black, inner sep=2.2pt] at (3,7.5) {};
  \node[font=\tiny\bfseries, text=red!70!black, anchor=south west] at (3.15,7.6)
    {Élargissement};

  % ◆ MVP Axe 1 livré (fin S4 — sem. 13)
  \node[diamond, fill=blue!65, draw=blue!80!black, inner sep=2.2pt] at (13,3.5) {};
  \node[font=\tiny\bfseries, text=blue!65!black, anchor=south west] at (13.15,3.6)
    {MVP Axe 1};

  % ◆ Axe 2 complet (fin S5 — sem. 17)
  \node[diamond, fill=green!55!black, draw=green!65!black, inner sep=2.2pt] at (17,2.5) {};
  \node[font=\tiny\bfseries, text=green!55!black, anchor=south west] at (17.15,2.6)
    {Axe 2 complet};

  % ◆ Soutenance (fin S7 — sem. 20)
  \node[diamond, fill=orange!80, draw=orange!75!black, inner sep=2.2pt] at (20,0.5) {};
  \node[font=\tiny\bfseries, text=orange!70!black, anchor=south east] at (19.9,0.6)
    {Soutenance};

  % ─── Légende jalons ─────────────────────────────────────────
  \node[font=\tiny, anchor=west] at (-8,-0.55) {%
    \textbf{Jalons :}
    \quad
    \tikz[baseline=-0.6ex]\node[diamond,fill=red!72,draw=red!80!black,inner sep=1.4pt]{};~Élargissement
    \quad
    \tikz[baseline=-0.6ex]\node[diamond,fill=blue!65,draw=blue!80!black,inner sep=1.4pt]{};~MVP Axe 1
    \quad
    \tikz[baseline=-0.6ex]\node[diamond,fill=green!55!black,draw=green!65!black,inner sep=1.4pt]{};~Axe 2 complet
    \quad
    \tikz[baseline=-0.6ex]\node[diamond,fill=orange!80,draw=orange!75!black,inner sep=1.4pt]{};~Soutenance
  };

\end{tikzpicture}
\end{figure}

\section{Bilan du chapitre}

Ce premier chapitre a posé le cadre institutionnel et stratégique du projet en
retraçant fidèlement sa trajectoire réelle : un PFE initialement centré sur la
modélisation stratégique, dont le périmètre a été élargi à la digitalisation après
la détection proactive d'un problème de gouvernance de la donnée interne. Les
six objectifs opérationnels, la méthodologie de travail agile et le planning en
huit sprints ont été présentés. Le chapitre suivant établit les fondements métier
et conceptuels nécessaires à la compréhension de ces deux axes.

\clearpage

% ============================================================
%  CHAPITRE 2 : FONDEMENTS MÉTIER, ANALYSE DES BESOINS ET CONCEPTION
% ============================================================

\end{document}
